\chapter{Sofware development rules}

\section{Algorithms}

\section{Flowchart}
    \subsection{Symbols}
        \url{https://en.wikipedia.org/wiki/Flowchart}
        \paragraph{Parallelogram}
            input/output

\section{Fundamental Principles}
    
    \subsection{Separation of Concerns (SoC)}
        Each part of the system should handle only one specific aspect (UI, logic, data, etc.).
    
    \subsection{Information Hiding}
        Implementation details should be hidden, and only an interface should be exposed.
    
    \subsection{Don’t Repeat Yourself (DRY)}
        No duplicated code; shared logic must exist in a single place.
    
    \subsection{Keep It Simple, Stupid (KISS)}
        Keep things simple; avoid unnecessary complexity.
    
    \subsection{You Aren’t Gonna Need It (YAGNI)}
        Don’t implement features you don’t currently need.
    
    \subsection{Law of Demeter (LoD) / Principle of Least Knowledge}
        A module should only interact with its immediate neighbors (“don’t talk to strangers”).

\section{Architecture Principles}
    
    \subsection{Layered Architecture Principle}
        The system should be divided into clear layers (presentation, application, domain, infrastructure).
    
    \subsection{Hexagonal Architecture (Ports and Adapters)}
        Dependencies should point inward to the core domain, and input/output should be handled through adapters.
    
    \subsection{CQRS (Command Query Responsibility Segregation)}
        Separate operations that change data (Commands) from those that read data (Queries).
    
    \subsection{Event-Driven Design Principle}
        Components should communicate via messages/events as much as possible, instead of direct dependencies.